% ----------------
% 节: 内容排版
% ----------------
\section{内容排版}
\subsection{字体}
\begin{frame}{Info.}
	\textbf{本小节将介绍字体相关设置, 包括字号、字体属性、字体强调、文本间距、文本修饰、标点符号与颜色.}
	\begin{multicols}{2}
		\begin{itemize}
			\item Establish: \textcolor{scugreen}{v1.0a (2021/11/30)}
			\item Update: \textcolor{scugreen}{v1.3a (2022/03/16)}
      \item Update: \textcolor{scugreen}{v1.3b (2022/04/13)}
      \item Update: \textcolor{scugreen}{v2.1c (2026/05/20)}
      \item Update: \textcolor{scugreen}{v2.1d (2026/06/15)}
		\end{itemize}
	\end{multicols}
	\mycopyright
\end{frame}

\begin{frame}{字号}
	\textbf{设置命令}: \\
	应用到环境全局: \cmd{begin}\marg{env} {\color{scured}\cmd{cmd}} \ldots \cmd{end}\marg{env}\\
	应用到特定文本: {\color{scured}\cmd{cmd}}\marg{text}\par\vspace{1ex}
	\begin{columns}[onlytextwidth]
		\begin{column}{.66\textwidth}
			\begin{table}[h]
				\centering
				\begin{tabular}{r|l}
				  \alert{命令}~\cmd{cmd} & \alert{效果} \\
					 			 \cmd{Huge} & \Huge{Fu-Gui老王}                        \\
						 		 \cmd{huge} & \huge{Fu-Gui老王}                        \\
							  \cmd{LARGE} & \LARGE{Fu-Gui老王}                       \\
							  \cmd{Large} & \Large{Fu-Gui老王}                       \\
							  \cmd{large} & \large{Fu-Gui老王}                       \\
					 \cmd{normalsize} & \normalsize{\alert{Fu-Gui老王}} (Default)\\
							  \cmd{small} & \small{Fu-Gui老王}                       \\
				 \cmd{footnotesize} & \footnotesize{Fu-Gui老王}                \\
					 \cmd{scriptsize} & \scriptsize{Fu-Gui老王}                  \\
								 \cmd{tiny} & \tiny{Fu-Gui老王}
				\end{tabular}
			\end{table}
		\end{column}
		\begin{column}{.25\textwidth}
			Beamer全局默认字号为 \alert{normalsize}, 一般情况不必修改默认字号.\\[3ex]
			\textbf{掌门大课堂}~字号局部调整建议使用左侧中 \cmd{footnotesize} 到 \cmd{Large} 间的字号.\\[3ex]

			老王, 形翼门自称超级英雄第二, 道上人称Fu-Gui叔, 其特异能力是形态的缩放.
		\end{column}
	\end{columns}
\end{frame}

\begin{frame}{字体属性}
	\textbf{掌门大课堂}~两类方法指定字体属性:
	(1) \alert{调用命令} \cmd{cmd}\marg{text} $\rightarrow$ 参数内文本, (2) \alert{声明命令} \{\cmd{cmd} \ldots\} $\rightarrow$ 环境全局文本.\vspace*{-1.1ex}
	\begin{table}[h]
		\centering
		\begin{tabular}{>{\bfseries\color{PrimaryC}}m{.05\textwidth}R{.26\textwidth}m{.16\textwidth}m{.12\textwidth}R{.26\textwidth}}
		  & \alert{\bfseries 名称} & \alert{\bfseries 调用命令} & \alert{\bfseries 声明命令} & \alert{\bfseries 中英文显示效果} \\[.4ex]
			\toprule
			\multirow{2}{=}{\centering 字体\\[-2pt]系列}
		  & 中等 Medium             & \cmd{textmd}\marg{text} & \cmd{mdseries} & Master Ma 形翼马掌门          \\
			& 加粗 Bold               & \cmd{textbf}\marg{text} & \cmd{bfseries} & \textbf{Master Ma 形翼马掌门} \\
			\cmidrule(lr){2-5}
			\multirow{4}{=}{\centering 字体\\[-2pt]形状}
			& 直立 Upright            & \cmd{textup}\marg{text} & \cmd{upshape}  & \textup{Master Ma 形翼马掌门} \\
			& 意大利斜体 Italic       & \cmd{textit}\marg{text} & \cmd{itshape}  & \textit{Master Ma 形翼马掌门} \\
			& 倾斜 Slanted            & \cmd{textsl}\marg{text} & \cmd{slshape}  & \textsl{Master Ma 形翼马掌门} \\
			& 小型大写 Small Caps     & \cmd{textsc}\marg{text} & \cmd{scshape}  & {\scshape Master Ma}    \\
			\cmidrule(lr){2-5}
			\multirow{3}{=}{\centering 字体\\[-2pt]族}
			& 衬线 Serif (罗马 Roman) & \cmd{textrm}\marg{text} & \cmd{rmfamily} & \textrm{Master Ma 形翼马掌门} \\
			& 无衬线 Sans-serif       & \cmd{textsf}\marg{text} & \cmd{sffamily} & \textsf{Master Ma 形翼马掌门} \\
			& 打字机 Typewriter       & \cmd{texttt}\marg{text} & \cmd{ttfamily} & \texttt{Master Ma 形翼马掌门} \\
			\bottomrule\\[-7pt]
			\multirow{5}{=}{\centering 字体\\[-2pt]属性\\[-1pt]组合}
			& Bold Italic       & \multicolumn{2}{c|}{{\bfseries\itshape Master Ma 形翼马掌门}} & \multirow{2}{=}{\centering\alert{\bfseries 常用场景}} \\
			& Bold Slanted      & \multicolumn{2}{c|}{{\bfseries\slshape Master Ma 形翼马掌门}} \\
			& Bold Small Caps   & \multicolumn{2}{c|}{{\bfseries\scshape Master Ma}}           & Bold $\rightarrow$ 强调重点            \\
			& Italic Small Caps & \multicolumn{2}{c|}{{\itshape\scshape Master Ma}}            & Italic $\rightarrow$ 注释/书名         \\
			& Italic Typewriter & \multicolumn{2}{c|}{{\itshape\ttfamily Master Ma 形翼马掌门}} & Typewriter $\rightarrow$ 代码/变量/URL \\
			\bottomrule
		\end{tabular}
	\end{table}
	\vspace*{-1.5ex}\textit{注: \textup{(1)}部分字体可能缺失Italic、Slanted或 Small Caps 字形, 此时 \LaTeX{} 会尝试使用近似字形或自动合成近似效果, 并在编译日志中给出警告. \textup{(2)}部分字体Italic Typewriter字形过于花哨, 建议类似本主题中使用Slanted Typewriter替代.}
\end{frame}

\begin{frame}{字体强调}
	老王常用三条口诀, 能迅速发挥出三种不同的狂暴力量, 一者健壮, 二者耀眼, 三者独特.\\

  \structure{结构字体}\\[1ex]
  \textbf{调用命令}: \cmd{structure}\textcolor{scugrey}{<\Arg{overlay}>}\marg{text}\hfill\structure{Neighbor 老王}\\
  \textbf{调用环境}:\\
  \cmd{begin}\marg*{structureenv}\textcolor{scugrey}{<\Arg{overlay}>}\\
  \hspace*{1em}\Arg{environment contents}\\
  \cmd{end}\marg*{structureenv}\\[1.5ex]

  \structure{重点字体}\\[1ex]
  \textbf{调用命令}: \cmd{alert}\textcolor{scugrey}{<\Arg{overlay}>}\marg{text}\hfill\alert{Neighbor 老王}\\
  \textbf{调用环境}:\\
  \cmd{begin}\marg*{alertenv}\textcolor{scugrey}{<\Arg{overlay}>}\\
  \hspace*{1em}\Arg{environment contents}\\
  \cmd{end}\marg*{alertenv}\\[1.5ex]

  \structure{强调字体}\\[1ex]
  \textbf{调用命令}: \cmd{emph}\marg{text}\hfill\emph{Neighbor 老王}\\
  标准\LaTeX{}强调命令, 中文模式下默认显示为\uline{下划线}(中文排版惯例), 英文模式下默认显示为\emph{斜体}, 会根据上下文自动切换.
\end{frame}

\begin{frame}{文本间距}
	\begin{columns}[T, onlytextwidth]
		\begin{column}{.7\textwidth}
			\structure{间距填充}
			\begin{table}[h]
				\centering
				\begin{tabular}{lll|lll}
					\alert{命令}~\cmd{cmd} & \alert{间距/功能} & \alert{效果} & \alert{命令}~\cmd{cmd} & \alert{间距(约)} & \alert{效果} \\
					\midrule
					\multicolumn{3}{c|}{\alert{\bfseries\scshape Text Spacing}}
					& \multicolumn{3}{c}{\alert{\bfseries\scshape Math Spacing\label{goto:MathQuad}}} \\
					\alert{\ttfamily\textbackslash\char32} & 普通空格 & a\ b
					& \cmd{,} & \texttt{3/18 em} & a\,b\,c \\
					\alert{\textasciitilde} & 不折行空格 & a~b
					& \cmd{thinspace} & 等效 \cmd{,} & a\,b\,c \\
					\cmd{enspace} & \texttt{1/2 em} & a\enspace b
					& \cmd{:} & \texttt{4/18 em} & a\:b\:c \\
					\cmd{quad} & \texttt{1 em} & a\quad b
					& \cmd{>} & \texttt{4/18 em} & a\:b\:c \\
					\cmd{qquad} & \texttt{2 em} & a\qquad b
					& \cmd{medspace} & 等效 \cmd{:} & a\:b\:c \\
					%
					\multicolumn{3}{c|}{\alert{\bfseries\scshape Glue}}
					& \cmd{;} & \texttt{5/18 em} & a\;b\;c \\
					\cmd{hfil} & \multicolumn{2}{l|}{水平弹性}
					& \cmd{thickspace} & 等效 \cmd{;} & a\;b\;c \\
					\cmd{hfill} & \multicolumn{2}{l|}{水平弹性(撑满)}
					& \cmd{!} & \texttt{-3/18 em} & a\!b~a\!b\!c \\
					\cmd{vfil} & \multicolumn{2}{l|}{垂直弹性}
					& \cmd{negthinspace} & 等效 \cmd{!} & a\!b~a\!b\!c \\
					\cmd{vfill} & \multicolumn{2}{l|}{垂直弹性(撑满)}
					& \cmd{negmedspace} & \texttt{-4/18 em} & a\negmedspace b~a\negmedspace b\negmedspace c \\
					%
					\multicolumn{3}{c|}{\alert{\bfseries\scshape Par Spacing}}
					& \cmd{negthickspace} & \texttt{-5/18 em} & a\negthickspace b~a\negthickspace b\negthickspace c \\
					\cmd{smallskip} & \multicolumn{2}{l|}{段落间小间距} & \scugoto{back:MathQuad}{BACK} \\
					\cmd{medskip}   & \multicolumn{2}{l|}{段落间中间距} & \\
					\cmd{bigskip}   & \multicolumn{2}{l|}{段落间大间距} & \\
				\end{tabular}
			\end{table}
		\end{column}
		\begin{column}{.28\textwidth}
			\structure{段落与换行}\\~\\
			\textbf{掌门大课堂}
			\begin{itemize}
				\item 段落结束直接换行: 空一行、\cmd{par}、\cmd{newline}
				\item 换行后使用预设长度控制段落间距: \cmd{smallskip}, \cmd{medskip}, \cmd{bigskip}
				\item 换行, 且可通过 \Arg{length} 控制段落间距: \cmd{\textbackslash}\oarg*{\Arg{length}}
				\item 换行但禁止分页: \cmd{\textbackslash*}\oarg*{\Arg{length}}
				\item 灵活小技巧(简单换两行, 句尾使用): \cmd{\textbackslash}\alert{\fverb{\textasciitilde}}\cmd{\textbackslash}
			\end{itemize}
			\vfill
		\end{column}
	\end{columns}
\end{frame}

\begin{frame}{文本修饰及标点符号}
	\begin{columns}[onlytextwidth]
		\begin{column}{.72\textwidth}
			\structure{文本修饰}\\[1ex]
			来自 \pkg{ulem} 宏包, 提供下划线类与删除线类修饰命令.
			\begin{table}[h]
				\centering
				\begin{tabular}{
					>{\centering\arraybackslash}m{.11\columnwidth}
					>{\centering\arraybackslash}m{.13\columnwidth}
					>{\centering\arraybackslash}m{.1\columnwidth}|
					>{\centering\arraybackslash}m{.11\columnwidth}
					>{\centering\arraybackslash}m{.13\columnwidth}
					>{\centering\arraybackslash}m{.1\columnwidth}
				}
					\multicolumn{3}{c}{\alert{\bfseries 下划线类(underline)}} & \multicolumn{3}{c}{\alert{\bfseries 删除线类(strikeout)}} \\
					\midrule
					下划线   & \cmd{uline}  & \uline{混元}  & 删除线   & \cmd{sout} & \sout{弹抖} \\
					双下划线 & \cmd{uuline} & \uuline{形翼} & 斜删除线 & \cmd{xout} & \xout{闪电鞭} \\
					波浪线   & \cmd{uwave}      & \uwave{太极门} & \\
					虚线	   & \cmd{dashuline}  & \dashuline{五连鞭} & \\
					点线	   & \cmd{dotuline}   & \dotuline{松果} & \\
				\end{tabular}
			\end{table}
		\end{column}
		\begin{column}{.27\textwidth}
			\textbf{简单自定义纹理修饰:}\\[0.5ex]
			\cmd{newcommand}\cmd{cmd}\marg*{\%\\
			  \hspace*{1em}\cmd{bgroup}\%\\
				\hspace*{1em}\cmd{markoverwith}\marg{pattern}\%\\
    		\hspace*{1em}\cmd{ULon}\%\\
			}\\[0.5ex]
			具体实现可参考 \pkg{ulem} 官方文档和 \href{https://tex.stackexchange.com/questions/530454}{\color{scublue}TeX.SE 社区相关讨论}.
		\end{column}
	\end{columns}\vspace*{-.5ex}
	\begin{columns}[onlytextwidth]
		\begin{column}{.25\textwidth}
			\structure{标点符号}\\[1ex]
			中西文混排可统一使用全角标点符号, \alert{自然科学类}或\alert{纯西文}排版建议使用\alert{半角标点符号}. \\[.5ex]
			\CTeX{} \cmd{punctstyle}\marg{style} 可在导言区设置标点样式. \Arg{style} 可选 \texttt{quanjiao} \alert{全角}、\texttt{banjiao} \alert{半角}、\texttt{kaiming} \alert{开明} (。！？等部分标点为全角).
		\end{column}
		\begin{column}{.74\textwidth}
			\hspace*{1.2em}\textbf{掌门大课堂}
			\begin{itemize}
				\item \alert{连字}: 部分西文字母连用时存在连字, 可通过 \alert{\fverb{shelf\{\}ful}} 或 \alert{\fverb{shelf\char`\\/ful}} 来取消.
				\item \alert{引号}: 中文单/双引号可直接输入. 西文单引号为 \alert{\fverb{`\,\textquotesingle}}, 双引号为 \alert{\fverb{``\,\textquotesingle\textquotesingle}}, \alert{\fverb{`}} 在ESC键下方.
				\item \alert{连字符}: 中文破折号在键盘键入\alert{一次全角长横线}(即两个连续的 \alert{\fverb{—}} )或 使用\alert{两次} \cmd{textemdash}. 西文连字符分单词连接 \alert{\fverb{-}}, 短连接符 \alert{\fverb{--}} 和破折号 \alert{\fverb{---}}.
				\item \alert{省略号}: 中文在键盘键入\alert{六次点号}或 使用\alert{两次} \cmd{ldots}. 西文使用\alert{三个点号} 或 \cmd{ldots}.
				\item \alert{特殊符号}: \LaTeX{} 保留字符 \alert{\texttt{@}} \alert{\fverb{\#}} \alert{\fverb{\$}} \alert{\fverb{\%}} \alert{\fverb{\&}} \alert{\fverb{\{}} \alert{\fverb{\}}} \alert{\fverb{\_}} 需在输入字符前加反斜杠 \alert{\fverb{\textbackslash}}. \alert{\textasciitilde} 和 \alert{\textasciicircum} 分别使用 \cmd{textasciitilde} 和 \cmd{textasciicircum}, 反斜杠使用 \cmd{textbackslash}.
			\end{itemize}
		\end{column}
	\end{columns}
\end{frame}

\begin{frame}{颜色}
	\begin{columns}[onlytextwidth]
		\begin{column}{.54\textwidth}
			\textbf{掌门大课堂}\\[1ex]
			定义新颜色: \cmd{definecolor}\marg{name}\marg{model}\marg{spec} \\[.7ex]
			基于已有颜色创建别名或混合颜色: \\
			\cmd{colorlet}\marg{name}\marg{existing} \\[.7ex]
			使用颜色着色文字: \cmd{textcolor}\marg{color}\marg{text} \\[.7ex]
			使用颜色着色文字: {\cmd{color}\marg{color} \Arg{text}} \\[.7ex]
			使用颜色着色背景: \cmd{colorbox}\marg{color}\marg{text} \\[.7ex]
			使用颜色着色并加边框: \\
			\cmd{fcolorbox}\marg{framecolor}\marg{backgroundcolor}\marg{text}
		\end{column}
		\begin{column}{.43\textwidth}
			\hfill\tiny\setlength{\tabcolsep}{0pt}
			\begin{tabular}{rr*{10}{l}}
				& & \rotatebox{70}{\texttt{100\%}} & \rotatebox{70}{\texttt{10\%}} &
				\rotatebox{70}{\texttt{20\%}} & \rotatebox{70}{\texttt{30\%}} &
				\rotatebox{70}{\texttt{40\%}} & \rotatebox{70}{\texttt{50\%}} &
				\rotatebox{70}{\texttt{60\%}} & \rotatebox{70}{\texttt{70\%}} &
				\rotatebox{70}{\texttt{80\%}} & \rotatebox{70}{\texttt{90\%}} \\
				\texttt{scured}\hspace*{1.5em}    & \textcolor{scured}{锦绣红}\hspace*{1.5em}    &
				\swatchC{scured}      & \swatchC{scured10}    & \swatchC{scured20}    &
				\swatchC{scured30}    & \swatchC{scured40}    & \swatchC{scured50}    &
				\swatchC{scured60}    & \swatchC{scured70}    & \swatchC{scured80}    &
				\swatchC{scured90}    \\
				\texttt{scugrey}\hspace*{1.5em}   & \textcolor{scugrey}{优雅灰}\hspace*{1.5em}   &
				\swatchC{scugrey}     & \swatchC{scugrey10}   & \swatchC{scugrey20}   &
				\swatchC{scugrey30}   & \swatchC{scugrey40}   & \swatchC{scugrey50}   &
				\swatchC{scugrey60}   & \swatchC{scugrey70}   & \swatchC{scugrey80}   &
				\swatchC{scugrey90}   \\
				\texttt{scublue}\hspace*{1.5em}   & \textcolor{scublue}{宝石蓝}\hspace*{1.5em}   &
				\swatchC{scublue}     & \swatchC{scublue10}   & \swatchC{scublue20}   &
				\swatchC{scublue30}   & \swatchC{scublue40}   & \swatchC{scublue50}   &
				\swatchC{scublue60}   & \swatchC{scublue70}   & \swatchC{scublue80}   &
				\swatchC{scublue90}   \\
				\texttt{scugreen}\hspace*{1.5em}  & \textcolor{scugreen}{荷叶绿}\hspace*{1.5em}  &
				\swatchC{scugreen}    & \swatchC{scugreen10}  & \swatchC{scugreen20}  &
				\swatchC{scugreen30}  & \swatchC{scugreen40}  & \swatchC{scugreen50}  &
				\swatchC{scugreen60}  & \swatchC{scugreen70}  & \swatchC{scugreen80}  &
				\swatchC{scugreen90}  \\
				\texttt{scuyellow}\hspace*{1.5em} & \textcolor{scuyellow}{银杏黄}\hspace*{1.5em} &
				\swatchC{scuyellow}   & \swatchC{scuyellow10} & \swatchC{scuyellow20} &
				\swatchC{scuyellow30} & \swatchC{scuyellow40} & \swatchC{scuyellow50} &
				\swatchC{scuyellow60} & \swatchC{scuyellow70} & \swatchC{scuyellow80} &
				\swatchC{scuyellow90} \\
			\end{tabular}\\[.5ex]
			\scriptsize
			本模板依据四川大学 VIS 手册定义了 \textcolor{scured}{\texttt{scured}} \textcolor{scugrey}{\texttt{scugrey}} \textcolor{scublue}{\texttt{scublue}} \textcolor{scugreen}{\texttt{scugreen}} \textcolor{scuyellow}{\texttt{scuyellow}} 5个颜色系列, 添加后缀 \texttt{10}--\texttt{90} , 如 \textcolor{scured70}{\texttt{scured70}} \textcolor{scublue80}{\texttt{scublue80}} 为各系列浅色变体.
		\end{column}
	\end{columns}
	
	\definecolor{exrgb}{RGB}{255,0,0}
	\definecolor{exhtml}{HTML}{FF0000}
	\definecolor{exgray}{gray}{0.5}
	\definecolor{excmyk}{cmyk}{.1, .2, .3, .4}
	\colorlet{expurple}{red!50!blue}
	\begin{table}[h]
		\centering
		\begin{tabular}{clclc}
			\alert{模式} & \alert{命令} & \alert{参数取值} & \alert{示例命令} & \alert{效果} \\
			\midrule
			rgb  & \cmd{definecolor} & \texttt{0}--\texttt{1} &
			\cmd{definecolor}\marg*{myred}\marg*{rgb}\marg*{1, 0, 0} & \colorbox{exrgb}{\phantom{XX}} \\
			RGB  & \cmd{definecolor} & \texttt{0}--\texttt{255} &
			\cmd{definecolor}\marg*{myred}\marg*{RGB}\marg*{255, 0, 0} & \colorbox{exrgb}{\phantom{XX}} \\
			HTML & \cmd{definecolor} & 6位十六进制 &
			\cmd{definecolor}\marg*{myred}\marg*{HTML}\marg*{FF0000} & \colorbox{exhtml}{\phantom{XX}}\\
			gray & \cmd{definecolor} & \texttt{0}--\texttt{1} &
			\cmd{definecolor}\marg*{mygray}\marg*{gray}\marg*{0.5} & \colorbox{exgray}{\phantom{XX}} \\
			cmyk & \cmd{definecolor} & \texttt{0}--\texttt{1} &
			\cmd{definecolor}\marg*{mygold}\marg*{cmyk}\marg*{.1, .2, .3, .4} & \colorbox{excmyk}{\phantom{XX}} \\
			混合 & \cmd{colorlet}    & \texttt{color!num!color} &
			\cmd{colorlet}\marg*{mypurple}\marg*{red!50!blue} & \colorbox{expurple}{\phantom{XX}}
		\end{tabular}
	\end{table}
	\vspace*{-1.5ex}
\end{frame}

\subsection{数学公式}
\begin{frame}{Info.}
	\textbf{本小节将介绍数学公式的排版, 包括行内公式、行间公式、多行公式、数学字体命令、数学符号、数学修饰、定界符与矩阵.}
	\begin{multicols}{2}
		\begin{itemize}
			\item Establish: \textcolor{scugreen}{v1.0a (2021/11/30)}
			\item Update: \textcolor{scugreen}{v1.3a (2022/03/16)}
      \item Update: \textcolor{scugreen}{v1.3b (2022/04/13)}
      \item Update: \textcolor{scugreen}{v2.1d (2026/06/15)}
		\end{itemize}
	\end{multicols}
	注: 自 \textcolor{scugreen}{v1.3a (2021/03/16)} 起, 本小节已自手册中\alert{\bfseries 数学环境}小节拆出.\par
	\mycopyright
\end{frame}

\begin{frame}{行内与行间公式}
	\textbf{掌门大课堂}~数学模式中普通空格会被忽略, 需使用专用空格命令\label{back:MathQuad} \scugoto{goto:MathQuad}{DTL}; 若需在公式中插入普通文本, 请使用 \cmd{text}\marg{text}.\\
	\begin{columns}[T, onlytextwidth]
		\begin{column}{.45\columnwidth}
			\structure{行内公式}\\[1ex]
			\textbf{使用方式}: \fverb{\$\Arg{equation}\$} 和 \fverb{\cmd{(}\Arg{equation}\cmd{)}}\\[1ex]
			\textbf{例如}, 最概然速率为 $v_p = \sqrt{\dfrac{2kT}{\mu}} = \sqrt{\dfrac{2RT}{M}}$, 其中$R$是气体常数, $M = N_A \mu$是物质的摩尔质量.\\[2ex]

			\structure{行间公式(无编号)}\\[1ex]
			\textbf{使用方式}: \fverb{\$\$\Arg{equation}\$\$} 和 \fverb{\cmd{[}\Arg{equation}\cmd{]}}\\[.5ex]
			和 \env{displaymath} 或 \env{equation*} 环境:\\[.5ex]
			\begin{columns}[onlytextwidth]
				\begin{column}{.5\textwidth}
					\cmd{begin}\marg*{displaymath}\\
					\hspace*{1em}\Arg{equation}\\
					\cmd{end}\marg*{displaymath}
				\end{column}
				\begin{column}{.5\textwidth}
					\cmd{begin}\marg*{equation*}\\
					\hspace*{1em}\Arg{equation}\\
					\cmd{end}\marg*{equation*}
				\end{column}
			\end{columns}\vspace*{1ex}
			\textbf{例如}, 麦克斯韦分布函数为,\\[-3ex]
			\begin{equation*}
				f(v) = \dfrac{\mathrm{d}N}{N\,\mathrm{d}v} = 4\pi \Big(\dfrac{\mu}{2\pi kT}\Big)^{3/2} v^2 \mathrm{exp}\Big(-\dfrac{\mu v^2}{2kT}\Big)
			\end{equation*}
		\end{column}
		\begin{column}{.5\columnwidth}
			\structure{行间公式(有编号)}\\[1ex]
			\textbf{使用方式}: \env{equation} 环境.\\[.5ex]
			\cmd{begin}\marg*{equation}\cmd{label}\marg{label}\\
			\hspace*{1em}\Arg{equation}\\
			\hspace*{1em}\cmd{tag}\marg{tag}\\
			\cmd{end}\marg*{equation}\\[.5ex]
			\Arg{label} 建议以 \texttt{eq:} 开头, \env{equation} 支持使用 \cmd{tag} 自定义编号格式 \\[1ex]
			\textbf{例如}, 平均速率为,\\[-1ex]
			\begin{equation}
				\bar{v} = \int_0^\infty vf(v)\,\mathrm{d}v = \sqrt{\dfrac{8kT}{\pi\mu}} = \sqrt{\dfrac{8RT}{\pi M}}
			\end{equation}
			\textbf{例如(自定义编号格式)}, 方均根速率为,\\[-3ex]
			\begin{equation}
				v_{rms} = \Big(\int_0^\infty v^2f(v)\,\mathrm{d}v\Big)^{1/2} = \sqrt{\dfrac{3kT}{\mu}} = \sqrt{\dfrac{3RT}{M}}
				\tag{$\alpha$}
			\end{equation}
		\end{column}
	\end{columns}
\end{frame}

\begin{frame}{多行公式(1/3)}
	\begin{columns}[onlytextwidth]
		\begin{column}{.36\columnwidth}
			\alert{\bfseries multline 环境}~
			用于长公式折行, 首行左对齐, 末行右对齐, 中间行居中, 整体产生一个编号.\\[1ex]
			\textbf{使用方式}:\\[.5ex]
			\cmd{begin}\marg*{multline}\\
			\hspace*{1em}\Arg{line 1}\cmd{\textbackslash}\\
			\hspace*{1em}\ldots\\
			\hspace*{1em}\Arg{line n}\cmd{label}\marg{label}\\
			\cmd{end}\marg*{multline}\\[3ex]

			\alert{\bfseries gather 环境}~
			多行公式各自居中, 每行单独编号, \env{gather*} 环境则不编号.\\[1ex]
			\textbf{使用方式}:\\[.5ex]
			\cmd{begin}\marg*{gather}\\
			\hspace*{1em}\Arg{equation 1}\cmd{label}\marg{label 1}\cmd{\textbackslash}\\
			\hspace*{1em}\Arg{equation 2}\cmd{label}\marg{label 2}\cmd{\textbackslash}\\
			\hspace*{1em}\ldots\\
			\cmd{end}\marg*{gather}
		\end{column}
		\begin{column}{.6\columnwidth}
			\textbf{例如(multline)}, Riemann 和展开计算,\\[-4ex]
			\begin{multline}
				A=\lim_{n\rightarrow\infty}\Delta x\left(a^{2}+\left(a^{2}+2a\Delta x+\left(\Delta x\right)^{2}\right)\right.\\
				+\left(a^{2}+2\cdot2a\Delta x+2^{2}\left(\Delta x\right)^{2}\right)\\
				+\left(a^{2}+2\cdot3a\Delta x+3^{2}\left(\Delta x\right)^{2}\right)\\
				+\ldots\\
				\left.+\left(a^{2}+2\cdot(n-1)a\Delta x+(n-1)^{2}\left(\Delta x\right)^{2}\right)\right)\\
				=\frac{1}{3}\left(b^{3}-a^{3}\right)
			\end{multline}
			\textbf{例如(gather)}, 勾股定理、欧拉恒等式与二项式定理,\\[-3ex]
			\begin{gather}
				a^2 + b^2 = c^2 \\
				e^{i\pi} + 1 = 0 \\
				(x+y)^n = \sum_{k=0}^{n} \binom{n}{k} x^{n-k} y^k
			\end{gather}
		\end{column}
	\end{columns}
\end{frame}

\begin{frame}{多行公式(2/3)}
	\begin{columns}[onlytextwidth]
		\begin{column}{.36\columnwidth}
			\alert{\bfseries align 环境}~
			多行对齐公式, 使用 \fverb{\&} 指定对齐位置, 可设置多个对齐列, 每行都编号. \cmd{notag} 或 \cmd{nonumber} 可取消某行编号, \env{align*} 则全部不编号.\\[1ex]
			\textbf{使用方式}:\\[.5ex]
			\cmd{begin}\marg*{align}\\
			\hspace*{1em}\Arg{expr 1} \fverb{\&} \Arg{expr 2}\cmd{label}\marg{label 1}\cmd{\textbackslash}\\
			\hspace*{1em}\Arg{expr 3} \fverb{\&} \Arg{expr 4}\cmd{notag}\cmd{\textbackslash} \\
			\cmd{end}\marg*{align}\\[2ex]

			\alert{\bfseries aligned 环境}~
			类似 \env{align}, 但需嵌套于 \env{equation} 等数学环境中, 整体产生一个编号.\\[1ex]
			\textbf{使用方式}:\\[.5ex]
			\cmd{begin}\marg*{equation}\cmd{label}\marg{label}\\
			\hspace*{1em}\cmd{begin}\marg*{aligned}\\
			\hspace*{2em}\Arg{expr 1} \fverb{\&} \Arg{expr 2} \cmd{\textbackslash}\\
			\hspace*{2em}\Arg{expr 3} \fverb{\&} \Arg{expr 4}\\
			\hspace*{1em}\cmd{end}\marg*{aligned}\\
			\cmd{end}\marg*{equation}
		\end{column}
		\begin{column}{.6\columnwidth}
			\textbf{例如(align)}, 质能方程,
			\begin{align}
				E=m&c^2 & &E=mc^2\\
				E=~&mc^2 & E&=mc^2 \notag \\
				E&=mc^2 & E=~&mc^2\\
				&E=mc^2 & E=m&c^2
			\end{align}
			\textbf{例如(aligned)}, 旋转参考系中的矢量求导,
			\begin{equation}
				\begin{aligned}
				\dfrac{\mathrm{d}}{\mathrm{d}t} \symbfit{f} &= \dfrac{\mathrm{d}f_x}{\mathrm{d}t} \symbfit{\hat{i}} + \dfrac{\mathrm{d}\symbfit{\hat{i}}}{\mathrm{d}t} f_x + \dfrac{\mathrm{d}f_y}{\mathrm{d}t} \symbfit{\hat{j}} + \dfrac{\mathrm{d}\symbfit{\hat{j}}}{\mathrm{d}t} f_y + \dfrac{\mathrm{d}f_z}{\mathrm{d}t} \symbfit{\hat{k}} + \dfrac{\mathrm{d}\symbfit{\hat{k}}}{\mathrm{d}t} f_z \\
				&= \dfrac{\mathrm{d}f_x}{\mathrm{d}t} \symbfit{\hat{i}} + \dfrac{\mathrm{d}f_y}{\mathrm{d}t} \symbfit{\hat{j}} + \dfrac{\mathrm{d}f_z}{\mathrm{d}t} \symbfit{\hat{k}} + \big[\symbf{\Omega} \times \big( f_x \symbfit{\hat{i}} + f_y \symbfit{\hat{j}} + f_z \symbfit{\hat{k}}\big) \big] \\
				&= \Big(\dfrac{\mathrm{d}\symbfit{f}}{\mathrm{d}t}\Big)_r + \symbf{\Omega} \times \symbfit{f}(t)
				\end{aligned}
			\end{equation}
		\end{column}
	\end{columns}
\end{frame}

\begin{frame}{多行公式(3/3)}
	\begin{columns}[onlytextwidth]
		\begin{column}{.49\columnwidth}
			\alert{\bfseries dcases 环境}~
			用于排版分段函数或条件定义, 采用 displaystyle 数学样式(区别于 \env{cases} 环境使用 textstyle 数学样式, 会导致数学公式缩小), 需调用 \pkg{mathtools} 宏包.\\[1ex]
			\textbf{使用方式}:\\[.5ex]
			\cmd{begin}\marg*{equation}\cmd{label}\marg{label}\\
			\hspace*{1em}\cmd{begin}\marg*{dcases}\\
			\hspace*{2em}\Arg{case 1} \cmd{\textbackslash}\\
			\hspace*{2em}\ldots\\
			\hspace*{1em}\cmd{end}\marg*{dcases}\\
			\cmd{end}\marg*{equation}\\[2ex]

			\alert{\bfseries split 环境}~
			用于单个公式的对齐与拆行, 需嵌套于 \env{equation} 等数学环境中, 整体产生一个编号, 仅支持单列对齐.\\[1ex]
			\textbf{使用方式}:\\[.5ex]
			\cmd{begin}\marg*{equation}\cmd{label}\marg{label}\\
			\hspace*{1em}\cmd{begin}\marg*{split}\\
			\hspace*{2em}\Arg{expr 1} \fverb{\&} \Arg{expr 2} \cmd{\textbackslash}\\
			\hspace*{2em}\Arg{expr 3} \fverb{\&} \Arg{expr 4}\\
			\hspace*{1em}\cmd{end}\marg*{split}\\
			\cmd{end}\marg*{equation}
		\end{column}\hspace*{.01\columnwidth}
		\begin{column}{.49\columnwidth}
			\textbf{例如(dcases)}, Maxwell 方程组,
			\begin{equation}
				\begin{dcases}
				\oint_l \symbfit{H} \cdot \mathrm{d}\symbfit{l} = \iint_S \symbfit{J} \cdot \mathrm{d}\symbfit{S} + \iint_S \dfrac{\partial\symbfit{D}}{\partial t} \cdot \mathrm{d}\symbfit{S} \\
				\oint_l \symbfit{E} \cdot \mathrm{d}\symbfit{l} = - \iint_S \dfrac{\partial\symbfit{B}}{\partial t} \cdot \mathrm{d}\symbfit{S} \\
				\oint_S \symbfit{B} \cdot \mathrm{d}\symbfit{S} = 0 \\
				\oint_S \symbfit{D} \cdot \mathrm{d}\symbfit{S} = \iiint_V \rho \mathrm{d}V
				\end{dcases}
			\end{equation}
			\textbf{例如(split)}, 变量代换求解积分,
			\begin{equation}
				\begin{split}
				\int_0^{\pi} x \sin x \,\mathrm{d}x &= -x \cos x \Big|_0^{\pi} + \int_0^{\pi} \cos x \,\mathrm{d}x \\
				&= \pi + \sin x \Big|_0^{\pi} \\
				&= \pi
				\end{split}
			\end{equation}\vspace*{-1ex}
		\end{column}
	\end{columns}
\end{frame}

\begin{frame}{数学字体命令}
	\textsc{Native} 为标准\LaTeX{}数学字体命令, \textsc{Ams} 为 \pkg{amsfonts} 或 \pkg{amssymb} 宏包提供的扩展数学字体命令. \textsc{Unicode} 为 \pkg{unicode-math} 宏包提供的数学字体命令 (详见 \pkg{unicode-math} 手册).
	\begin{table}[h]
		\centering\setlength{\tabcolsep}{2.5pt}
		\begin{tabular}{>{\bfseries\color{PrimaryC}}m{.08\textwidth}
			R{.15\textwidth}R{.2\textwidth}m{.06\textwidth}
			R{.15\textwidth}R{.2\textwidth}m{.06\textwidth}}
			& \alert{\bfseries 名称} & \alert{\bfseries 命令} & \alert{\bfseries 效果}
			& \alert{\bfseries 名称} & \alert{\bfseries 命令} & \alert{\bfseries 效果} \\
			\midrule
			\multirow{5}{=}{\centering \textsc{Native}}
		  & 正常数学样式 & \cmd{mathnormal}\marg{text} & $\mathnormal{Rr}$ \\
			& 直立        & \cmd{mathrm}\marg{text}     & $\mathrm{Rr}$     &
			  花体        & \cmd{mathcal}\marg{text}    & $\mathcal{Rr}$    \\
			& 粗体        & \cmd{mathbf}\marg{text}     & $\mathbf{Rr}$     &
			  无衬线      & \cmd{mathsf}\marg{text}     & $\mathsf{Rr}$     \\
			& 斜体        & \cmd{mathit}\marg{text}     & $\mathit{Rr}$     &
			  打字机      & \cmd{mathtt}\marg{text}     & $\mathtt{Rr}$     \\
			\cmidrule(lr){2-7}
			\multirow{1}{=}{\centering \textsc{Ams}}
		  & 黑板粗体    & \cmd{mathbb}\marg{text}     & $\mathbb{Rr}$     &
			  哥特体      & \cmd{mathfrak}\marg{text}   & $\mathfrak{Rr}$   \\
			\cmidrule(lr){2-7}
			\multirow{10}{=}{\centering \textsc{Unicode}}
			& 正常数学样式 & \cmd{symnormal}\marg{text}  & $\symnormal{Rr}$  \\
		  & 直立        & \cmd{symup}\marg{text}      & $\symup{Rr}$      &
			  粗直立      & \cmd{symbfup}\marg{text}    & $\symbfup{Rr}$    \\
			& 斜体        & \cmd{symit}\marg{text}      & $\symit{Rr}$      &
			  粗斜体      & \cmd{symbfit}\marg{text}    & $\symbfit{Rr}$    \\
			& 花体        & \cmd{symcal}\marg{text}     & $\symcal{Rr}$     &
			  粗花体      & \cmd{symbfcal}\marg{text}   & $\symbfcal{Rr}$   \\
			& 无衬线直立   & \cmd{symsfup}\marg{text}    & $\symsfup{Rr}$    &
			  粗无衬线直立 & \cmd{symbfsfup}\marg{text}  & $\symbfsfup{Rr}$  \\
			& 无衬线斜体   & \cmd{symsfit}\marg{text}    & $\symsfit{Rr}$    &
			  粗无衬线斜体 & \cmd{symbfsfit}\marg{text}  & $\symbfsfit{Rr}$  \\
			& 手写体      & \cmd{symscr}\marg{text}     & $\symscr{Rr}$     &
			  粗手写体    & \cmd{symbfscr}\marg{text}   & $\symbfscr{Rr}$   \\
			& 哥特体      & \cmd{symfrak}\marg{text}    & $\symfrak{Rr}$    &
			  粗哥特体    & \cmd{symbffrak}\marg{text}  & $\symbffrak{Rr}$  \\
			& 黑板粗体    & \cmd{symbb}\marg{text}      & $\symbb{Rr}$      &
			  黑板意大利体 & \cmd{symbbit}\marg{text}    & $\symbbit{Rr}$    \\
			& 打字机      & \cmd{symtt}\marg{text}      & $\symtt{Rr}$      &
				粗无衬线    & \cmd{symbfsf}\marg{text}    & $\symbfsf{Rr}$    \\
			\midrule
		\end{tabular}
	\end{table}
\end{frame}

\begin{frame}{数学符号}
	\vspace*{-6pt}
	\begin{table}[h]
		\centering\setlength{\tabcolsep}{1.5pt}\renewcommand{\arraystretch}{.8}
		\begin{tabular}{>{\bfseries\color{PrimaryC}}m{.04\textwidth}
			m{.13\textwidth}m{.04\textwidth}
			m{.134\textwidth}m{.04\textwidth}
			m{.13\textwidth}m{.04\textwidth}
			m{.13\textwidth}m{.04\textwidth}
			m{.13\textwidth}m{.04\textwidth}}
			\midrule
			\multirow{5}{=}{\centering \\[-8pt]小写\\[-1.8pt]希腊}
		  & \quad\cmd{alpha}      & $\alpha$      &
			  \quad\cmd{beta}       & $\beta$       &
			  \quad\cmd{gamma}      & $\gamma$      &
			  \quad\cmd{delta}      & $\delta$      &
			  \quad\cmd{epsilon}    & $\epsilon$    \\
			& \quad\cmd{zeta}       & $\zeta$       &
			  \quad\cmd{eta}        & $\eta$        &
			  \quad\cmd{theta}      & $\theta$      &
			  \quad\cmd{iota}       & $\iota$       &
			  \quad\cmd{lambda}     & $\lambda$     \\
			& \quad\cmd{mu}         & $\mu$         &
			  \quad\cmd{nu}         & $\nu$         &
			  \quad\cmd{xi}         & $\xi$         &
			  \quad\cmd{pi}         & $\pi$         &
			  \quad\cmd{rho}        & $\rho$        \\
			& \quad\cmd{sigma}      & $\sigma$      &
			  \quad\cmd{tau}        & $\tau$        &
			  \quad\cmd{phi}        & $\phi$        &
			  \quad\cmd{chi}        & $\chi$        &
			  \quad\cmd{psi}        & $\psi$        \\
			& \quad\cmd{omega}      & $\omega$      &
			  \quad\cmd{varepsilon} & $\varepsilon$ &
			  \quad\cmd{varphi}     & $\varphi$     &
			  \quad\cmd{varpi}      & $\varpi$      & \\
			\cmidrule(lr){2-11}
			\multirow{2}{=}{\centering \vspace*{1pt}大写\\[-4pt]希腊\vspace*{-2pt}}
		  & \quad\cmd{Gamma}      & $\Gamma$      &
			  \quad\cmd{Delta}      & $\Delta$      &
			  \quad\cmd{Theta}      & $\Theta$      &
			  \quad\cmd{Lambda}     & $\Lambda$     &
			  \quad\cmd{Xi}         & $\Xi$         \\
			& \quad\cmd{Pi}         & $\Pi$         &
			  \quad\cmd{Sigma}      & $\Sigma$      &
			  \quad\cmd{Phi}        & $\Phi$        &
			  \quad\cmd{Psi}        & $\Psi$        &
			  \quad\cmd{Omega}      & $\Omega$      \\
			\midrule
			\multirow{3}{=}{\centering \\[-9pt]数学\\[-1.8pt]函数}
			& \quad\cmd{sin}				& $\sin$        &
			  \quad\cmd{cos}        & $\cos$        &
			  \quad\cmd{tan}        & $\tan$        &
			  \quad\cmd{log}        & $\log$        &
			  \quad\cmd{ln}         & $\ln$         \\
			& \quad\cmd{exp}        & $\exp$        &
			  \quad\cmd{lim}        & $\lim$        &
				\quad\cmd{max}        & $\max$        &
				\quad\cmd{min}        & $\min$        &
			  \quad\cmd{det}        & $\det$        \\
			& \quad\cmd{dim}        & $\dim$        &
			  \quad\cmd{ker}        & $\ker$        &
			  \quad\cmd{deg}        & $\deg$        &
				\quad\cmd{bmod}       & $\bmod$       \\
			\midrule
		\end{tabular}
	\end{table}\vspace*{-23.2pt}
	\begin{table}[h]
		\centering\setlength{\tabcolsep}{1.5pt}\renewcommand{\arraystretch}{.8}
		\begin{tabular}{>{\bfseries\color{PrimaryC}}m{.04\textwidth}
			R{.09\textwidth}m{.08\textwidth}m{.04\textwidth}
			R{.09\textwidth}m{.08\textwidth}m{.04\textwidth}
			R{.09\textwidth}m{.08\textwidth}m{.04\textwidth}
			R{.09\textwidth}m{.08\textwidth}m{.04\textwidth}}
			\midrule
			\multirow{3}{=}{\centering \\[-8pt]运算\\[-1.8pt]符号}
		  & 乘       & \cmd{times}    & $\times$    &
			  除       & \cmd{div}      & $\div$      &
			  点乘     & \cmd{cdot}     & $\cdot$     &
			  加减     & \cmd{pm}       & $\pm$       \\
			& 减加     & \cmd{mp}       & $\mp$       &
			  星乘     & \cmd{ast}      & $\ast$      &
			  并集     & \cmd{cup}      & $\cup$      &
			  交集     & \cmd{cap}      & $\cap$      \\
			& 异或     & \cmd{oplus}    & $\oplus$    & \\
			\cmidrule(lr){2-13}
			\multirow{3}{=}{\centering \\[-7pt]关系\\[-1.8pt]符号}
		  & 约等于   & \cmd{approx}   & $\approx$     &
			  全等     & \cmd{equiv}    & $\equiv$      &
			  不等于   & \cmd{neq}      & $\neq$        &
			  小于等于 & \cmd{leq}      & $\leq$        \\
			& 大于等于 & \cmd{geq}      & $\geq$        &
			  远小于   & \cmd{ll}       & $\ll$         &
			  远大于   & \cmd{gg}       & $\gg$         &
			  约等于   & \cmd{simeq}    & $\simeq$      \\
			& 相似于   & \cmd{sim}      & $\sim$        &
			  成比例   & \cmd{propto}   & $\propto$     &
			  属于     & \cmd{in}       & $\in$         &
			  子集     & \cmd{subset}   & $\subset$     \\
			\midrule
		\end{tabular}
	\end{table}\vspace*{-23.2pt}
	\begin{table}[h]
		\centering\setlength{\tabcolsep}{2pt}\renewcommand{\arraystretch}{.8}
		\begin{tabular}{>{\bfseries\color{PrimaryC}}m{.04\textwidth}
			R{.09\textwidth}m{.151\textwidth}m{.04\textwidth}
			R{.09\textwidth}m{.1455\textwidth}m{.04\textwidth}
			R{.09\textwidth}m{.151\textwidth}m{.04\textwidth}}
			\midrule
			\multirow{3}{=}{\centering \\[-8pt]箭头}
		  & 右箭头   & \cmd{rightarrow}  & $\rightarrow$  &
			  左箭头   & \cmd{leftarrow}   & $\leftarrow$   &
			  双向     & \cmd{leftrightarrow} & $\leftrightarrow$ \\
			& 长右箭头 & \cmd{longrightarrow} & $\longrightarrow$ &
			  长左箭头 & \cmd{longleftarrow}  & $\longleftarrow$  &
			  推出     & \cmd{Rightarrow}  & $\Rightarrow$  \\
			& 等价     & \cmd{Leftrightarrow} & $\Leftrightarrow$ &
			  映射     & \cmd{mapsto}      & $\mapsto$      &
			  蕴含     & \cmd{Leftarrow}   & $\Leftarrow$   \\
			\midrule
		\end{tabular}
	\end{table}
	\vspace*{-6pt}
\end{frame}

\begin{frame}{数学修饰}
	\begin{table}[h]
		\centering\setlength{\tabcolsep}{3pt}
		\begin{tabular}{>{\bfseries\color{PrimaryC}}m{.04\textwidth}
			R{.1\textwidth}R{.26\textwidth}m{.06\textwidth}
			R{.1\textwidth}R{.26\textwidth}m{.06\textwidth}}
			& \alert{\bfseries 名称} & \alert{\bfseries 命令} & \alert{\bfseries 效果}
			& \alert{\bfseries 名称} & \alert{\bfseries 命令} & \alert{\bfseries 效果} \\
			\midrule
			\multirow{4}{=}{\centering 重音}
		  & 帽子     & \cmd{hat}\marg{content}       & $\hat{a}$         &
			  宽帽子   & \cmd{widehat}\marg{content}   & $\widehat{adf}$   \\
			& 波浪号   & \cmd{tilde}\marg{content}     & $\tilde{a}$       &
			  宽波浪号 & \cmd{widetilde}\marg{content} & $\widetilde{adf}$ \\
			& 短横线   & \cmd{bar}\marg{content}       & $\bar{a}$         &
			  向量箭头 & \cmd{vec}\marg{content}       & $\vec{a}$         \\
			& 点号     & \cmd{dot}\marg{content}       & $\dot{a}$         &
			  双点号   & \cmd{ddot}\marg{content}      & $\ddot{a}$        \\
			\cmidrule(lr){2-7}
			\multirow{3}{=}{\centering 上方\\[-2pt]修饰}
		  & 上划线   & \cmd{overline}\marg{content}          & $\overline{adf}$          &
			  上方括号 & \cmd{overbrace}\marg{content}         & $\overbrace{adf}$         \\
			& 上方堆叠 & \cmd{overset}\marg{top}\marg{content} & $\overset{\text{def}}{=}$ &
			  右箭头   & \cmd{overrightarrow}\marg{content}    & $\overrightarrow{adf}$    \\
			& 左箭头   & \cmd{overleftarrow}\marg{content}     & $\overleftarrow{adf}$     \\
			\cmidrule(lr){2-7}
			\multirow{3}{=}{\centering 下方\\[-2pt]修饰}
		  & 下划线   & \cmd{underline}\marg{content}           & $\underline{adf}$    &
			  下方括号 & \cmd{underbrace}\marg{content}          & $\underbrace{adf}$   \\
			& 下方堆叠 & \cmd{underset}\marg{down}\marg{content} & $\underset{\sim}{a}$ &
			  右箭头   & \cmd{xrightarrow}\marg{content}         & $\xrightarrow{adf}$  \\
			& 左箭头   & \cmd{xleftarrow}\marg{content}          & $\xleftarrow{adf}$   \\
			\cmidrule(lr){2-7}
			\multirow{2}{=}{\centering 删除\\[-2pt]线}
		  & 斜删除线   & \cmd{cancel}\marg{content}   & $\cancel{adf}$  &
			  反斜删除线 & \cmd{bcancel}\marg{content}  & $\bcancel{adf}$ \\
			& 交叉删除线 & \cmd{xcancel}\marg{content}  & $\xcancel{adf}$ \\
			\midrule
		\end{tabular}
	\end{table}
	\vspace*{-2ex}~~\textit{注: 删除线命令 \textup{\cmd{cancel} \cmd{bcancel} \cmd{xcancel}} 需调用 \textup{\pkg{cancel}} 宏包.}
\end{frame}

\begin{frame}{定界符与矩阵}
	\begin{columns}[T, onlytextwidth]
		\begin{column}{.5\columnwidth}
			\structure{定界符}\\[1ex]
			固定高度, 成对使用:\\\vspace*{-2ex}
			\begin{multicols}{2}
				\fverb{(}\Arg{content}\fverb{)} \hfill $(x)$\\[.5ex]
				\fverb{[}\Arg{content}\fverb{]} \hfill $[x]$\\[.5ex]
				\cmd{\{}\Arg{content}\cmd{\}} \hfill $\{x\}$\\[.5ex]
				\fverb{|}\Arg{content}\fverb{|} \hfill $|x|$\\
			\end{multicols}\vspace*{-2ex}
			\cmd{langle}\Arg{content}\cmd{rangle} \hfill $\langle x \rangle$\\[2ex]
			自动匹配高度, 成对使用:\\[.5ex]
			\cmd{left}\Arg{symbol}\Arg{content}\cmd{right}\Arg{symbol}\\[.5ex]
			如 \Arg{symbol} 为圆括号, \cmd{left(}\ldots\cmd{right)}; 若某侧无需显示定界符, 可使用 \fverb{.} 占位, 如 \cmd{left.}\ldots\cmd{right|}; 若公式中有竖线, 可以使用 \cmd{middle|} 自动调整中间竖线高度.
			\vspace*{-1ex}\[
				\left(\frac{a^2+b^2}{c^2+d^2}\right) \quad
				\left[\sum_{i=1}^{n} x_i\right] \quad
				\left\{x \mid x>0\right\} \quad
				\left\langle \frac{\partial f}{\partial x} \middle| \frac{\partial f}{\partial y} \right\rangle
			\]\vspace*{-1ex}
			固定尺寸:\\[.5ex]
			\cmd{cmd}\Arg{symbol}\Arg{content}\cmd{cmd}\Arg{symbol} \hfill $\bigl( \Bigl( \biggl( \Biggl( x \Biggr) \biggr) \Bigr) \bigr)$\\[-1.5ex]
			\cmd{cmd} 可为 \cmd{big}, \cmd{Big}, \cmd{bigg}, \cmd{Bigg}
		\end{column}
		\begin{column}{.48\columnwidth}
			\structure{矩阵}\\[1ex]
			\pkg{amsmath} 宏包给出了6种常用的矩阵环境\\[.5ex]
			无定界符: \env{matrix}\\[.5ex]
			有定界符: \\\vspace*{-2ex}
			\begin{multicols}{2}
				\env{pmatrix} \hfill $\big(\cdots\big)$\\[.5ex]
				\env{bmatrix} \hfill $\big[\cdots\big]$\\[.5ex]
				\env{Bmatrix} \hfill $\bigl\{\cdots\bigr\}$\\[.5ex]
				\env{vmatrix} \hfill $|\cdots|$\\[.5ex]
				\env{Vmatrix} \hfill $||\cdots||$\\~\\
			\end{multicols}\vspace*{-2ex}
			\LaTeX{} 也提供了 类似 \env{tabular} 的 \env{array} 环境, 支持自定义列对齐(\fverb{l}, \fverb{c}, \fverb{r}), 也可配合 \cmd{left}/\cmd{right} 添加定界符.\\[1ex]
			\textbf{使用方法}: \\[1ex]
			\begin{columns}[onlytextwidth]
				\begin{column}{.5\textwidth}
					\cmd{begin}\marg*{bmatrix}\\
					\hspace*{1em}\Arg{content}\\
					\cmd{end}\marg*{bmatrix}
				\end{column}
				\begin{column}{.5\textwidth}
					\cmd{begin}\marg*{array}\marg{spec}\\
					\hspace*{1em}\Arg{content}\\
					\cmd{end}\marg*{array}
				\end{column}
			\end{columns}\vspace*{1ex}
			\textbf{例如},\vspace*{-1ex}
			\[
			\begin{bmatrix}
				E & 0 \\
				0 & E \\
			\end{bmatrix}
			\begin{pmatrix}
				E & 0 \\
				0 & E \\
			\end{pmatrix}
			\quad
			\left[\begin{array}{ccc}
			a_{11} & a_{12} & a_{13} \\
			a_{21} & a_{22} & a_{23}
			\end{array}\right]
			\]
		\end{column}
	\end{columns}
\end{frame}

%% End of file `basic-content.tex'.
